%! Planning a Study — Unit 3, Lesson 3.1 — AP Practice
% GRADUAL RELEASE: AP-format practice, extra credit, scored 1 to 5;
% the graded practice is the `homework` component that follows it.
% Two multiple-choice items (MC1: pick the experiment for a reading program — the distractors
% are careful observational designs; MC2: the appropriate generalization from a random sample
% of Ohio adults — only to the population sampled), then one free-response set on energy drinks
% and sleep: (a) observational or experiment, (b) does the study support "energy drinks cause
% insomnia" (no — no treatment imposed; a third variable could drive both), (c) the SPIRAL —
% name the display for two quantitative variables and the pattern to expect (Unit 2).
% Every rubric box is full width on its own line under its part, hung with \gradedhere over
% \writelines{8} (about 2.1in, room for the drawn 2in box). Each \gradedblank line, display=
% included, is BYTE-IDENTICAL in ap_practice_key/; key mode prints the display. The
% next-lesson preview is NOT here: it closes the packet, in ../homework/.
\documentclass[12pt]{article}
\usepackage{apstats-article}
\usepackage{apstats-boxes}
\usepackage{grading}
\gradingassignment{apstats-u03-l01-ap}
\gradingcourse{advanced_statistics}

\begin{document}

\pageheader{Unit 3, Lesson 3.1}{AP Practice}
{\small\itshape Write each final answer in its box. Only what is inside a box is graded.}\par
% NAMESTRIP: no \namedateperiod here — the name row lives on the cover only.

\vspace{0.06in}

% Authored BYTE-IDENTICALLY in ap_practice/ and ap_practice_key/.
\boxguard[8]
\begin{remindbox}
\small
This page is \textbf{extra credit, scored 1 to 5} --- your graded homework is the last section
of this packet. Work these AP-format reps for practice and bring questions to study hall.
\end{remindbox}

\vspace{0.06in}

\boxguard[10]
\begin{headlinebox}{goldbg}
\small
\textbf{Section I --- Multiple Choice.} Write the letter of the single best answer in its Letter box.
\end{headlinebox}

\vspace{0.06in}

\begin{enumerate}[leftmargin=1.9em, itemsep=0.10in]

  \item A school district wants to know whether a new reading program \emph{causes} higher
        reading scores for second graders. Which plan is an experiment?
        \begin{enumerate}[label=\textbf{(\Alph*)}, leftmargin=2.2em, itemsep=1pt]
          \item Compare the scores of students whose parents enrolled them in the program with
                the scores of students who did not enroll.
          \item Survey second-grade teachers about whether the program seems to work.
          \item Randomly assign 60 second graders to the new program or the current one, then
                compare their scores.
          \item Record the reading scores of every second grader in the district for a year.
          \item Review records from schools that already chose to use the program.
                \hfill\gradedright{\textit{Letter:}~\gradedblank[id=ap1, type=choice, points=1, answer=C]{0.45in}{0.24in}}\rule[-0.2in]{0pt}{0.48in}
        \end{enumerate}

  \item A polling firm selects 500 adults at random from a list of all adults in Ohio. Of
        them, 62\% say they support a new state park. Which conclusion is most appropriate?
        \begin{enumerate}[label=\textbf{(\Alph*)}, leftmargin=2.2em, itemsep=1pt]
          \item About 62\% of all adults in the United States support the park.
          \item About 62\% of all adults in Ohio support the park, give or take some
                sampling variability.
          \item Exactly 62\% of all adults in Ohio support the park.
          \item The result describes only the 500 adults who were asked.
          \item Living in Ohio causes people to support new parks.
                \hfill\gradedright{\textit{Letter:}~\gradedblank[id=ap2, type=choice, points=1, answer=B]{0.45in}{0.24in}}\rule[-0.2in]{0pt}{0.48in}
        \end{enumerate}

\end{enumerate}

\vspace{0.08in}

\boxguard[27]
\begin{headlinebox}{goldbg}
\small
\textbf{Section II --- Free Response.} Every answer must be written \emph{in the context of the
problem} --- that is what earns credit on the AP exam.
\end{headlinebox}

\vspace{0.06in}

\begin{enumerate}[leftmargin=1.9em, itemsep=0.10in, resume]

  \item A health researcher reviews survey records for 1{,}500 adults in one county. Adults who
        drink two or more energy drinks a day report sleeping fewer hours per night, on average,
        than adults who drink none.

        \vspace{6pt}

        \boxguard[16]
        \textbf{(a)} Is this an observational study or an experiment? Explain.\\[4pt]
        \gradedhide{\gradedhere{\gradedblank[id=ap3a, type=rubric, points=2, display={Observational study. The researcher only looked at records of what the adults already did; no one was assigned to drink or not drink energy drinks.}]{\linewidth+6pt}{\gradedlinesheight{2}}}\writelines{8}}

        \vspace{4pt}
        \boxguard[16]
        \textbf{(b)} A headline reads, \emph{``Energy drinks cause sleepless nights.''} Does this
        study support the headline? Explain in context.\\[4pt]
        \gradedhide{\gradedhere{\gradedblank[id=ap3b, type=rubric, points=2, display={No. Without an imposed treatment the study shows only an association. Something else, such as night-shift work or stress, could lead people both to drink more energy drinks and to sleep less.}]{\linewidth+6pt}{\gradedlinesheight{2}}}\writelines{8}}

        \vspace{4pt}
        \boxguard[16]
        \textbf{(c)} Name a graph that displays the relationship between energy drinks per day and
        hours of sleep, and describe the pattern you would expect.\\[4pt]
        \gradedhide{\gradedhere{\gradedblank[id=ap3c, type=rubric, points=2, display={A scatterplot (drinks per day on the x-axis, hours of sleep on the y-axis). Expect a negative association: adults who drink more energy drinks tend to sleep fewer hours.}]{\linewidth+6pt}{\gradedlinesheight{2}}}\writelines{8}}

\end{enumerate}

\end{document}
